\chapter*{Abstract}
\addcontentsline{toc}{chapter}{ABSTRACT} % to add abstract to table of contents
\begin{spacing}{1.3}
Post-disaster search and rescue operations across catastrophic environments such as landslides, floods, and structural collapses face severe logistical constraints where the initial twenty-four to seventy-two hours determine victim survival. Conventional relief missions rely on ground search teams or manually piloted single drones, suffering from battery limits, telemetry bottlenecks, and high rescuer risk. To overcome these challenges, this project introduces Search And Rescue Autonomous Swarm (SARAS), an end-to-end coordinated multi-unmanned aerial vehicle framework for automated disaster zone exploration, edge victim detection, and agentic triage support. The proposed methodology establishes a dual-tier intelligence pipeline comprising onboard edge perception and decentralized multi-agent coordination. For onboard perception, an 8-bit integer quantized You Only Look Once version 8 nano detector is optimized for embedded companion computers. For multi-vehicle coordination, an asynchronous market auction mechanism dynamically partitions search sectors and executes rapid reassignment upon vehicle attrition. Evaluated across the VisDrone aerial benchmark and high-fidelity Gazebo multi-robot disaster simulations, the onboard vision model achieves an inference throughput of 16.8 frames per second with an 81.6\% mean average precision at 0.5 intersection-over-union, consuming under 6.5 watts. In multi-drone swarm trials, the decentralized protocol achieves 100\% spatial search coverage across a one-square-kilometer disaster theater in 11.4 minutes with a 96.6\% path efficiency, detecting vehicle failures within 6.0 seconds and reallocating orphaned sectors in under 420 milliseconds. Furthermore, an agentic triage system incorporates retrieval-augmented generation to provide incident commanders with protocol-compliant extraction guidelines. This work demonstrates a scalable, resilient aerial swarm platform bridging physical edge robotics with intelligent tactical decision support.
\end{spacing}
